% ====================================================================
\section{Introduction}\label{sec:intro}
% ====================================================================

\subsection{The problem}
A finite set $A\subset\R^d$ is \emph{rainbow} if its $\binom{|A|}{2}$ pairwise distances are pairwise distinct
(sets with at most two points are rainbow). A set $X$ is \emph{rainbow-$n$-free} if it contains no rainbow
subset of size $n$. For $n\ge 3$ and $d\ge 1$ let $f_d(n)$ be the least $m$ such that every set of $m$ points of
$\R^d$ contains a rainbow $n$-subset, and let
\[
  s_n(j):=\sup\{|X| : X\subset\R^{N}\text{ finite, rainbow-$n$-free, } \dim\aff X\le j\}
\]
(the ambient dimension $N$ is irrelevant). Then $f_d(n)=s_n(d)+1$, and $s_n(0)=1$.

Erd\H{o}s asked whether $f_d(n)=2^{o(d)}$ for every fixed $n$; this is Problem~\#1088 on the
Erd\H{o}s problems website~\cite{EP1088}.\footnote{Within the project the problem is also tracked under
the identifier JSP-000905.} We prove that the answer is yes, with a much stronger quantitative bound.

\begin{theorem}[Main Theorem]\label{thm:main-intro}
For every fixed $n\ge 3$ there is a constant $C_n$ such that every set of more than
$2^{C_n\sqrt d\,\log_2(d+2)}$ points of $\R^d$ contains $n$ points with all $\binom n2$ pairwise distances
distinct. That is, $f_d(n)=2^{O_n(\sqrt d\log d)}=2^{o(d)}$.
\end{theorem}

\subsection{History}
As recorded on the problem page~\cite{EP1088}, Erd\H{o}s and Straus showed $f_d(n)\le c_n^{\,d}$ for a constant
$c_n$; the problem asks whether the exponential dependence on $d$ is necessary. In the other direction,
$s$-distance sets with $s=\binom n2-1$ are rainbow-$n$-free, since a rainbow $n$-set needs $\binom n2$ distinct
distances. For instance the $0/1$-vectors of weight $w$ in $\R^{d+1}$ lie in a $d$-dimensional affine subspace
and determine only the distances $\sqrt{2i}$, $1\le i\le w$; with $w=\binom n2-1$ this gives
$f_d(n)>\binom{d+1}{\binom n2-1}$ for $d+1\ge \binom n2 -1$. So $f_d(n)$ grows at least polynomially in $d$, with
degree depending on $n$, and we are not aware of any super-polynomial lower bound.

For $n=3$ the problem is classical and polynomial: a set is rainbow-$3$-free exactly when every three of its
points span an isosceles triangle, where collinear triples with two equal distances count as (degenerate)
isosceles triangles. Blokhuis~\cite{Blokhuis} proved that an isosceles set in $\R^d$ has at most
$\binom{d+2}{2}$ points. His argument is purely metric, so it covers degenerate triples, and it gives
$f_d(3)\le\binom{d+2}2+1$. For $n\ge 4$ we know of no bound of the form $2^{o(d)}$ in the
literature before the present work.

\subsection{The new tool: Theorem LR}
The geometric part of our argument repeatedly produces matrices of the form $M[x,w]=|x-w|^2$, of rank at most
$D+2$, in which every entry is ``explained'' by one of boundedly many simple causes. A typical cause is
$|x-w|=|w-t|$ for a fixed point $t$, which depends only on the column $w$. Another is $|x-w|=|x-t|$, which
depends only on the row $x$. A third is $|x-w|=d$ for a fixed number $d$. The tool that converts such a cover
into geometric structure is the following rank--rectangle theorem (Theorem~\ref{thm:LR}).

\medskip\noindent\textbf{Theorem LR (informal).} \emph{Let $M\in\R^{m\times n}$ have rank at most $r$, and let
$\tau_1,\dots,\tau_k$ be functions on $[m]\times[n]$, each depending only on the row or only on the column. If every
entry satisfies $M[x,w]=\tau_i(x,w)$ for some $i$, then there are $R\subseteq[m]$, $C\subseteq[n]$ and a single
$i$ with $M=\tau_i$ on $R\times C$, $|R|\ge 2^{-L}m$ and $|C|\ge 2^{-L}n$, where
$L=O(k^3\sqrt{r+2}\,\log(r+2))$.}

\medskip
For Boolean matrices this is the rectangle form of Lovett's bound~\cite{Lovett}, in Rothvoss's
formulation~\cite{Rothvoss}. Our proof follows Rothvoss's hyperplane-rounding argument and adds two ingredients.
The first is a normal form, obtained on a sub-rectangle of constant relative size, which allows us to rescale the
matrix at its largest gap \emph{by a row function, a column function or a constant}. This keeps a margin even
when the type values accumulate. The second is a cleaning step, the Gavinsky--Lovett lemma (Lemma~\ref{lem:GL}),
applied to a Hadamard product, whose rank enters only logarithmically.

\begin{remark}[possible independent interest; to be confirmed by experts]\label{rem:SS}
The special case in which all types are column functions concerns \emph{$k$-listable} matrices, in which every
column takes at most $k$ distinct values. Here Theorem~LR yields a column-constant submatrix of relative size
$2^{-O(k^3\sqrt d\log d)}$ in each direction for rank-$d$ matrices (Corollary~\ref{cor:RQlist}). By
Singer--Sudan~\cite[Prop.~2.3]{SingerSudan}, a parallel $k$-partitioned point--hyperplane configuration in $\R^d$
yields a $k$-listable matrix of rank at most $d$. A column-constant (``$1$-listable'') submatrix of it then yields a
complete bipartite subgraph of the incidence graph, of proportional size up to a factor depending on $k$. Singer and
Sudan write that ``the best bound we know for parallel 3-partitioned configurations is $2^{-\tilde O(d)}$''
(\cite{SingerSudan}, arXiv version~5, Section~1). As far as we can tell, Corollary~\ref{cor:RQlist} improves this to
$2^{-O(\sqrt d\log d)}$ for every fixed $k$. Our literature search (summarised in Section~\ref{sec:concl}) was
limited, and the observation follows from Rothvoss's proof in a few lines, so experts may regard it as known. The
claim of improvement must be confirmed by experts before it is relied on. Theorem~LR does not bear on the
log-rank conjecture, which would require $2^{-\mathrm{polylog}(d)}$.
\end{remark}

\subsection{Prior in-project work}
This manuscript builds on earlier notes of a peer session in the same project~\cite{Peer}. Those notes
recorded the problem, its literature and the lower-bound constructions above. They also contained
polynomial-method ``peeling'' lemmas: a point all of whose distances to the rest of a set lie in a short list can
be removed using a low-degree polynomial in the lifted coordinates $(|y|^2,y)$. Lemma~H below
(Lemma~\ref{lem:H}) is a two-sided variant of those peeling lemmas, in which each pair needs to be ``killed'' from
one side only, at the price of squaring the dimension. The present argument does not use the other results of
those notes, which include reductions to prime sets and twin analyses.

\subsection{Verification status}
We state this plainly, because the result would resolve a problem of Erd\H{o}s and the argument is long.
\begin{itemize}[leftmargin=2em]
\item The mathematics was developed by AI agents (Claude, Anthropic) in sessions directed by the team.
\item Each component was checked by at least one independent AI referee pass. These passes found and repaired
  one genuine gap, a threshold gap in the top-level argument (fix F1, incorporated in
  Section~\ref{sec:top}), together with a number of expository gaps, all addressed here.
\item \textbf{Machine-checked.} The complete argument has been formalized in Lean~4 (Section~\ref{sec:lean}).
  The formal proof establishes the statement of Problem~\#1088 in the formal-conjectures repository verbatim. It
  depends only on the axioms \texttt{propext}, \texttt{Classical.choice} and \texttt{Quot.sound}, and it has no
  \texttt{sorry}. The formalization proves Theorem~LR with a weaker dependence on~$k$ than stated here; this suffices
  for the Main Theorem.
\item \textbf{No human expert has reviewed the informal argument.} Theorem~LR is new and unpublished, and every later
  step depends on it. It is the component we most want checked by specialists.
\item This manuscript consolidates the project notes. A few routine changes made in the consolidation, listed in
  Section~\ref{sec:ai}, have not themselves been refereed; among them are some explicit constants and the direct
  single-cap write-up of Theorem~R$^+$.
\item The finite, threshold-free implications were also tested by computer in exact arithmetic. Because every
  threshold in the argument is astronomically large, such tests are sanity checks only and cannot test the
  asymptotic claims (Section~\ref{sec:concl}).
\end{itemize}
Until an expert check has been carried out, the Main Theorem should be regarded as a \emph{candidate} solution.

\subsection{Organisation}
Section~\ref{sec:structure} gives an overview of the proof, the linear order of dependencies and the threshold
tower. Section~\ref{sec:prelim} collects the classical tools: John's theorem, Sheppard's formula, the
cleaning lemma (with proof), Lemma~H, the Caro--Wei and K\"onig theorems, and the orthogonal split. Section~\ref{sec:LR}
proves Theorem~LR. Section~\ref{sec:G1} proves the label lemma and its consequence~G1 for pairs with disjoint label
sets. Section~\ref{sec:framework} sets up the level-configurations, the hereditary relative bound~HR and the kill
step. Section~\ref{sec:top} treats the top three levels. Sections~\ref{sec:core} and~\ref{sec:slicing} prove the
class-count statement at the lower levels: the Core Lemma and the Slicing Lemma. Section~\ref{sec:main} closes the
induction and proves the Main Theorem. Section~\ref{sec:concl} contains concluding remarks and the computational
sanity checks, and Section~\ref{sec:ai} the statement on AI assistance and the acknowledgements.
